% !TeX program = lualatex
% !TeX encoding = utf8
% !TeX spellcheck = uk_UA
% !TeX root = ../main.tex

%% --------------------------------------------------------
%% Chapter file: Collisions.tex
%% --------------------------------------------------------

\graphicspath{{\currfilebase/Pictures}}

\Opensolutionfile{answer}[\currfilebase/\currfilebase-Answers]
\Writetofile{answer}{\protect\section*{Розділ~\ref*{sec:plasma-collisions}. \nameref*{sec:plasma-collisions}}}
\chapter{Зіткнення, провідність та перенос у~плазмі}\label{sec:plasma-collisions}\hypertarget{plasma-collisions}{}

Досі ми розглядали колективне екранування (розд.~\ref{sec:plasma-basics}) та рух окремих частинок у~заданих полях
(розд.~\ref{sec:particle-motion}), де окремі акти розсіяння були неважливі. Тепер додамо зіткнення. Це дасть три нові
явища: опір електричному струму, теплоперенос і~дифузію. Важливо відрізняти два рівні
опису: \emphx{колективність} пояснює реакцію плазми як цілого, а~\emphx{зіткнення}
визначають дисипацію імпульсу та енергії.

%% --------------------------------------------------------
\section{Кулонівські зіткнення}\index{Зіткнення!кулонівські|textbf}
%% --------------------------------------------------------

\begin{Theory}[Нагадаємо з~молекулярної фізики]
У~газі частинка швидкості $v_T$ із~перерізом зіткнення $\sigma$ у~середовищі концентрації $n$ має
час між зіткненнями $\tau=1/(nv_T\sigma)$ і~довжину вільного пробігу $\ell=v_T\tau$.
Рух між зіткненнями~--- випадкове блукання: після $N=t/\tau$ кроків довжиною $\ell$ маємо
$\langle x^2\rangle=N\ell^2=(\ell^2/\tau)\,t$, тобто дифузію з~коефіцієнтом $D\sim\ell^2/\tau\sim v_T^2\tau$
(точніше, $D=\tfrac13\ell v_T$). За тією самою логікою за порядком величини
в'язкість $\eta\sim nT\tau$, а~теплопровідність $\kappa\sim nT\tau/m$ (з~числовим множником порядку одиниці).
Далі всі оцінки виду «$D\sim v_T^2/\nu$, $\kappa\sim nv_T^2\tau$» ґрунтуються саме на цих співвідношеннях,
де $\nu=1/\tau$ і~$\lambda_{\mathrm{пр}}=\ell=v_T/\nu$.
\end{Theory}

Хоча колективні взаємодії в~плазмі домінують, зіткнення все ж відбуваються, і~саме вони
визначають опір плазми електричному струму, вирівнювання температур і~загасання хвиль.
Зіткнення заряджених частинок мають суттєву особливість порівняно зі зіткненнями
нейтральних атомів: сила далекодійна, тому повний імпульс частинка втрачає не в~одному
лобовому зіткненні, а~у~вигляді \emph{багатьох малих відхилень} від далеких партнерів.

\begin{figure}[ht]\centering
	\localinput{elpath.tikz}
	\caption{Шлях частинки в~нейтральному газі (різкі зіткнення) і~в~плазмі (плавне викривлення далекими прольотами).
	Довжину вільного пробігу $\ell$ означають як шлях, на якому напрям швидкості повертається на~$\sim\pi/2$}
	\label{fig:elpath}
\end{figure}

Електрон зі швидкістю $v$, що пролітає повз іон із~прицільною відстанню $b$, відхиляється
(формула Резерфорда) на кут $\theta$, де $\operatorname{tg}(\theta/2)=b_{90}/b$, а~
% ~~~~~~~~~~~~~~~~~~~~~~~~~~~~~~~~~~~~~~~~~~~~~~~~~~~~~~~~
\begin{equation*}
	b_{90}=\frac{e^2}{m_ev^2}\sim\frac{e^2}{T}
\end{equation*}
% ~~~~~~~~~~~~~~~~~~~~~~~~~~~~~~~~~~~~~~~~~~~~~~~~~~~~~~~~
(для $v\sim v_{Te}$)~--- відстань, при якій відхилення порівнянне з~прямим кутом. Для
$b\gg b_{90}$ відхилення мале, $\theta\approx2b_{90}/b$. Число проходжень з~прицільною
відстанню в~інтервалі $db$ пропорційне $b\,db$, а~внесок у~середній квадрат кута~---
$\theta^2\propto b^{-2}$, отже сумарний ефект усіх інтервалів пропорційний $\int db/b$:
кожна декада прицільних відстаней дає однаковий внесок. Інтегрування йде від $b_{90}$ (при
менших $b$ вже велике відхилення) до $\lambda_D$ (далі іон екранований). Звідси
\emphx{кулонівський логарифм}\index{Кулонівський логарифм|textbf}
% ~~~~~~~~~~~~~~~~~~~~~~~~~~~~~~~~~~~~~~~~~~~~~~~~~~~~~~~~
\begin{equation*}
	\ln\Lambda=\ln\frac{\lambda_D}{b_{90}}\sim\ln N_D\approx10\text{--}20.
\end{equation*}
% ~~~~~~~~~~~~~~~~~~~~~~~~~~~~~~~~~~~~~~~~~~~~~~~~~~~~~~~~
Він слабко залежить від параметрів плазми, і~для оцінок його вважають сталим.

\begin{Attention}
	Формула $\ln\Lambda\approx\ln N_D$ має сенс лише при $N_D=n\lambda_D^3\gg1$. Для сильнозв'язаної плазми $\ln\Lambda\lesssim1$, і~всі результати цього розділу втрачають точність.
\end{Attention}

Значення $\ln\Lambda$ для типових плазм (формула NRL: $\ln\Lambda=23-\ln\bigl(n^{1/2}T^{-3/2}\bigr)$ при $T<10$~еВ, $24-\ln\bigl(n^{1/2}T^{-1}\bigr)$ при $T>10$~еВ; $n$ у~$\mathrm{см}^{-3}$, $T$ в~еВ) наведено в~табл.~\ref{tab:coulomb-log}:
для дуже різних плазм $\ln\Lambda$ лежить у~межах $10$--$20$.

\begin{table}[ht]
	\caption{Кулонівський логарифм для типових плазм.}
	\label{tab:coulomb-log}
	\begin{platbl}{colspec={X[2,l,m]X[c,m]X[c,m]X[c,m]}}
		Плазма & $n$, $\mathrm{см}^{-3}$ & $T$, еВ & $\ln\Lambda$ \\
		Лабораторний розряд & $10^{11}$ & $1$ & $10{,}3$ \\
		Іоносфера & $10^{5}$ & $0{,}1$ & $13{,}8$ \\
		Токамак (край) & $10^{13}$ & $10^{2}$ & $13{,}6$ \\
		Сонячна корона & $10^{8}$ & $10^{2}$ & $19{,}4$ \\
		Токамак (серцевина) & $10^{14}$ & $10^{4}$ & $17{,}1$ \\
	\end{platbl}
\end{table}

Оцінимо частоту зіткнень «на пальцях». Одне далеке прольотне зіткнення з~прицільною відстанню $b$
повертає швидкість на кут $\theta\approx2b_{90}/b$. Кути від різних зіткнень випадкові, тому додаються
\emph{квадрати}; на шляху $dx$ електрон зустрічає $2\pi b\,n\,db\,dx$ розсіювачів, тож
$$\frac{d\langle\theta^2\rangle}{dx}=\int2\pi b\,n\,db\left(\frac{2b_{90}}b\right)^2
=8\pi nb_{90}^2\ln\frac{b_{\max}}{b_{\min}}=8\pi nb_{90}^2\ln\Lambda.$$
Умова $\langle\theta^2\rangle\sim1$ визначає довжину пробігу $\ell\sim1/(8\pi nb_{90}^2\ln\Lambda)$, звідки
$$\nu=\frac v\ell\approx\frac{8\pi ne^4\ln\Lambda}{m^2v^3}.$$
Звідси одразу видно: (1)~\emph{логарифм} з'являється тому, що кожна декада $b$ дає однаковий внесок;
(2)~далекі зіткнення домінують над близькими в~$\ln\Lambda$ разів; (3)~$\nu\propto v^{-3}$, тобто
$\nu\propto T^{-3/2}$ (звідси ефект Драйсера, див. нижче). Точний кінетичний розрахунок з~усередненням за
максвеллівським розподілом і~іншим означенням $v$ замінює $8\pi$ на~$4\sqrt{2\pi}/3$ (множник порядку одиниці)
і~дає частоту електрон-іонних зіткнень (для $Z=1$)
% ~~~~~~~~~~~~~~~~~~~~~~~~~~~~~~~~~~~~~~~~~~~~~~~~~~~~~~~~
\begin{equation}\label{eq:nu-ei}
	\nu_{ei}=\frac{4\sqrt{2\pi}}3\,\frac{n_ee^4\ln\Lambda}{m_e^{1/2}(T_e)^{3/2}}\approx2{,}9\times10^{-6}\,n_e[\mathrm{см}^{-3}]\,\ln\Lambda\,T_e^{-3/2}[\text{еВ}]~\mathrm{с}^{-1}.
\end{equation}
% ~~~~~~~~~~~~~~~~~~~~~~~~~~~~~~~~~~~~~~~~~~~~~~~~~~~~~~~~~~~~~~~~
Залежність $\nu\propto T^{-3/2}$ має просте пояснення: переріз кулонівського розсіяння
$\propto v^{-4}$, тому швидкі частинки \enquote{не помічають} партнерів. Для іонів
із~зарядом $Ze$ ($n_i=n_e/Z$) переріз розсіяння пропорційний $Z^2$, а~концентрація
розсіювачів~--- $1/Z$, тому $\nu_{ei}\propto Z$, а~провідність приблизно $\propto1/Z$.

Порівняємо частоту зіткнень із~плазмовою. Підстановка~\eqref{eq:nu-ei},
\eqref{eq:plasma-frequency} з~урахуванням~\eqref{eq:ND-Gamma} дає
% ~~~~~~~~~~~~~~~~~~~~~~~~~~~~~~~~~~~~~~~~~~~~~~~~~~~~~~~~
\begin{equation}\label{eq:collisionless}
	\tcbhighmath{\frac{\nu_{ei}}{\omega_{pe}}=\frac{2\sqrt2}3\ln\Lambda\,\Gamma^{3/2}\approx0{,}09\,\frac{\ln\Lambda}{N_D}}.
\end{equation}
% ~~~~~~~~~~~~~~~~~~~~~~~~~~~~~~~~~~~~~~~~~~~~~~~~~~~~~~~~~~~~~~~~
Це важливий, але тонкий результат: умова $N_D\gg1$ означає слабкий кулонівський зв'язок
і~велику кількість частинок у~дебаївській сфері, але сама по собі ще не визначає, чи є
конкретний процес беззіткнювальним. Довжина вільного пробігу $\lambda_{\mathrm{пр}}\sim
v_{Te}/\nu_{ei}$ справді набагато більша за $\lambda_D$ для типових слабкозв'язаних плазм,
але порівнювати її треба насамперед із~розміром системи $L$ або з~іншим характерним
масштабом задачі. Наприклад, для $n_e=10^{12}~\mathrm{см}^{-3}$, $T_e=10$~еВ,
$\ln\Lambda=15$ маємо $\nu_{ei}\approx1{,}4\times10^6~\mathrm{с}^{-1}$,
$\omega_{pe}\approx5{,}6\times10^{10}~\mathrm{с}^{-1}$,
$\nu_{ei}/\omega_{pe}\approx2{,}4\times10^{-5}$.

Для електрон-електронних зіткнень $\nu_{ee}\sim\nu_{ei}$, а~для іонів
$\nu_{ii}\sim(m_e/m_i)^{1/2}\nu_{ee}$ (іони важчі й~повільніші). Обмін енергією між
електронами та іонами ще повільніший, $\sim(m_e/m_i)\nu_{ei}$. Порівнюючи часи, для водню
маємо $\tau_{ee}:\tau_{ii}:\tau_{ei}^{\varepsilon}\sim1:43:1836$. Тому кожен компонент
спершу встановлює власний максвеллівський розподіл і~лише потім вирівнюється температура
між ними~--- це й~пояснює можливість $T_e\neq T_i$. Фізично: за одне зіткнення електрона з~іоном передається частка енергії $\sim m/M$,
тож для вирівнювання температур потрібно $\sim M/m$ зіткнень. Двотемпературна плазма $T_e\neq T_i$~--- звичайний стан.
Контрінтуїтивний факт: довжина вільного пробігу $\ell=v_T/\nu\propto T^2/(ne^4\ln\Lambda)$ \emph{не залежить від маси}, тому $\ell_e=\ell_i$ при $T_e=T_i$.



%% --------------------------------------------------------
\section{Провідність плазми}\index{Провідність плазми|textbf}
%% --------------------------------------------------------



Урахуємо зіткнення в~найпростішому наближенні, ввівши силу тертя об~іони. Рівняння руху
електрона в~полі $\vect E$
% ~~~~~~~~~~~~~~~~~~~~~~~~~~~~~~~~~~~~~~~~~~~~~~~~~~~~~~~~
\begin{equation*}
	m_e\frac{d\vect v}{dt}=-e\vect E-m_e\nu\vect v,
\end{equation*}
% ~~~~~~~~~~~~~~~~~~~~~~~~~~~~~~~~~~~~~~~~~~~~~~~~~~~~~~~~
де $\nu\equiv\nu_{ei}=1/\tau_e$. У~сталому полі швидкість встановлюється $\vect
v_d=-e\vect E/(m_e\nu)$, густина струму $\vect j=-n_ee\vect v_d=\sigma_0\vect E$
з~провідністю
% ~~~~~~~~~~~~~~~~~~~~~~~~~~~~~~~~~~~~~~~~~~~~~~~~~~~~~~~~
\begin{equation}\label{eq:drude}
	\tcbhighmath{\sigma_0=\frac{n_ee^2}{m_e\nu}=\frac{n_ee^2\tau_e}{m_e}}.
\end{equation}
% ~~~~~~~~~~~~~~~~~~~~~~~~~~~~~~~~~~~~~~~~~~~~~~~~~~~~~~~~~
У~$\nu$ входять електрон-іонні зіткнення. Зіткнення електронів між собою не змінюють їхній
сумарний імпульс, тому прямого тертя не створюють (вони лише впливають на форму розподілу,
тобто на числовий коефіцієнт). Для гармонічного поля $\vect E\propto e^{-i\omega t}$
отримуємо $\sigma(\omega)=\sigma_0/(1-i\omega/\nu)$, а~діелектрична проникність
$\varepsilon=1+4\pi i\sigma/\omega$ набуває вигляду
% ~~~~~~~~~~~~~~~~~~~~~~~~~~~~~~~~~~~~~~~~~~~~~~~~~~~~~~~~
\begin{equation}\label{eq:eps-collisions}
	\varepsilon(\omega)=1-\frac{\omega_p^2}{\omega(\omega+i\nu)}.
\end{equation}
% ~~~~~~~~~~~~~~~~~~~~~~~~~~~~~~~~~~~~~~~~~~~~~~~~~~~~~~~~~~~~~~
Дві граничні ситуації: при $\omega\gg\nu$ повертаємося до~\eqref{eq:plasma-permittivity},
а~при $\omega\ll\nu$ плазма поводиться як звичайний провідник ($\varepsilon\approx4\pi
i\sigma_0/\omega$). Уявна частина $\varepsilon$ описує поглинання. Середня за період
питома потужність джоулевих втрат $q=\tfrac12\operatorname{Re}\sigma\,E_0^2$,
$\operatorname{Re}\sigma=\sigma_0/(1+\omega^2/\nu^2)$. Саме так зіткнення в~нижній
іоносфері поглинають радіохвилі (розд.~\ref{sec:plasma-waves}), і~саме джоулеві втрати
спричиняють $R<1$ при відбитті від реальної плазми.



\subsection*{Повністю іонізована плазма}



Підставимо~\eqref{eq:nu-ei} у~\eqref{eq:drude}: концентрація \emph{скорочується}, і
% ~~~~~~~~~~~~~~~~~~~~~~~~~~~~~~~~~~~~~~~~~~~~~~~~~~~~~~~~
\begin{equation}\label{eq:spitzer}
	\sigma\propto\frac{T_e^{3/2}}{\ln\Lambda}.
\end{equation}
% ~~~~~~~~~~~~~~~~~~~~~~~~~~~~~~~~~~~~~~~~~~~~~~~~~~~~~~~~~
Точний кінетичний розрахунок (Спітцер) змінює числовий коефіцієнт приблизно вдвічі (для
струму вздовж поля $\sigma_\parallel\approx1{,}96\,n_ee^2/m_e\nu_{ei}$), і~для водневої
плазми
% ~~~~~~~~~~~~~~~~~~~~~~~~~~~~~~~~~~~~~~~~~~~~~~~~~~~~~~~~
\begin{equation*}
	\sigma_\parallel\approx1{,}7\times10^{14}\,\frac{T_e^{3/2}[\text{еВ}]}{\ln\Lambda}~\mathrm{с}^{-1}.
\end{equation*}
% ~~~~~~~~~~~~~~~~~~~~~~~~~~~~~~~~~~~~~~~~~~~~~~~~~~~~~~~~~
У~магнітному полі провідність стає анізотропною, але слід розрізняти два випадки.
\emph{Повністю іонізована плазма:} у~системі відліку плазми (з~урахуванням руху іонів) спітцерівський
опір поперек поля лише приблизно вдвічі більший, $\eta_\perp\approx2\eta_\parallel$, тобто магнітне
поле \emph{не} пригнічує поперечну провідність на порядки: електрон-іонне тертя зберігає сумарний
імпульс, а~дрейф $\vect E\times\vect B$ однаковий для обох сортів і~струму не створює.
\emph{Слабко іонізований газ} (іоносфера) або нерухомі іони: струм поперек поля підтримується лише
зіткненнями, і~педерсенівська провідність пригнічується фактором
$1/[1+(\omega_c/\nu)^2]$, так що за $\omega_c\gg\nu$ маємо
$\sigma_\perp/\sigma_\parallel\sim(\nu/\omega_c)^2\ll1$.

Висновок~\eqref{eq:spitzer} контрінтуїтивний: чим \emph{гарячіша} плазма, тим \emph{краще}
вона проводить. Причина в~тому, що зі зростанням температури електрони рухаються швидше,
але ефективність кулонівського розсіяння спадає ще швидше. Для $T_e=1$~кеВ
і~$\ln\Lambda=15$ маємо $\sigma\approx3{,}6\times10^{17}~\mathrm{с}^{-1}$, що вже
порівнянно з~провідністю міді при кімнатній температурі
($\sim5\times10^{17}~\mathrm{с}^{-1}$). Практичний наслідок: питома потужність омічного
нагрівання $j^2/\sigma\propto T^{-3/2}$ падає з~температурою, і~самим лише омічним
нагріванням плазму не можна довести до термоядерних температур (потрібне додаткове
нагрівання).

У~слабко іонізованому газі частоту зіткнень $\nu$ визначають зіткнення з~нейтральними атомами,
$\nu_{en}=n_n\sigma_{en}v_{Te}$, і~провідність~\eqref{eq:drude} залежить передусім від
концентрації електронів $n_e$, яка (через формулу Саха) швидко зростає з~температурою.

\subsection*{Поле Драйсера та «втікаючі» електрони}\label{sec:dreicer}

Оскільки $\nu_{ei}\propto v^{-3}$ ($\sigma\propto v^{-4}$), сила тертя на швидкий електрон
$F_{\rm fr}\sim m_e\nu_{ei}(v)\,v\propto v^{-2}$ \emph{спадає} зі~швидкістю. Тому у~сталому полі
$E$ електрон, який випадково набрав достатньої швидкості, вже не гальмується тертям, а~безперервно
прискорюється~--- стає \emph{втікаючим}\index{Втікаючі електрони|textbf} (runaway). Критичне поле, за якого сила $eE$ зрівнюється з~тертям на
тепловий електрон ($v\sim v_{Te}$), називають \emphx{полем Драйсера}\index{Драйсера поле|textbf}:
% ~~~~~~~~~~~~~~~~~~~~~~~~~~~~~~~~~~~~~~~~~~~~~~~~~~~~~~~~
\begin{equation}\label{eq:dreicer}
	\tcbhighmath{E_D=\frac{4\pi e^3n\ln\Lambda}{T_e}\approx2{,}6\times10^{-13}\,\frac{n\ln\Lambda}{T_e[\text{еВ}]}\ \text{В/см}.}
\end{equation}
% ~~~~~~~~~~~~~~~~~~~~~~~~~~~~~~~~~~~~~~~~~~~~~~~~~~~~~~~~
При $E<E_D$ електрони з~$v>v_c=v_{Te}\sqrt{E_D/E}$ все ж втікають, але їхня частка експоненційно мала
$\sim\exp(-v_c^2/2v_{Te}^2)=\exp(-E_D/2E)$. Для токамака ($n=10^{14}$~см$^{-3}$, $T_e=10$~кеВ):
$E_D\approx4$~В/м, тоді як електричне поле в~розряді $E\sim0{,}1$~В/м, і~втікаючих електронів практично немає.
Але під час швидкого зриву струму $T_e$ падає, а~$E$ зростає на порядки, виникають пучки електронів
з~енергіями в~десятки МеВ, небезпечні для стінок установки.

%% --------------------------------------------------------
\section{Теплопровідність плазми}\index{Теплопровідність плазми|textbf}
%% --------------------------------------------------------

Ті самі електрони переносять і~тепло. За порядком величини (кінетика газу) $\kappa_e\sim
n_ev_{Te}\lambda_{\mathrm{пр}}=n_ev_{Te}^2\tau_e$, де
$\tau_e=1/\nu_{ei}$. Оскільки $\tau_e\propto T_e^{3/2}/n_e$,
% ~~~~~~~~~~~~~~~~~~~~~~~~~~~~~~~~~~~~~~~~~~~~~~~~~~~~~~~~
\begin{equation*}
	\kappa_e\sim\frac{n_eT_e\tau_e}{m_e}\propto T_e^{5/2},
\end{equation*}
% ~~~~~~~~~~~~~~~~~~~~~~~~~~~~~~~~~~~~~~~~~~~~~~~~~~~~~~~~
і~вона, як і~провідність, не залежить від концентрації. Порівняння з~\eqref{eq:drude} дає
масштабне співвідношення $\kappa_e/\sigma\sim T_e/e^2$. Це аналог
співвідношення Відемана--Франца за температурною залежністю, але числовий коефіцієнт
визначається кінетичною моделлю і~не збігається з~металевим. Отже, гаряча плазма має дуже
високу теплопровідність: при нагріванні від $10$ до $100$~еВ вона зростає
у~$10^{5/2}\approx300$ разів. Теплопровідність іонів приблизно в~$\sqrt{m_e/m_i}$ разів
менша: $\kappa\propto\tau/m$, а~$\tau_i/\tau_e\sim\sqrt{m_i/m_e}$.

Це стосується переносу вздовж поля. У~замагніченій плазмі довжина кроку дорівнює не
$\lambda_{\mathrm{пр}}$, а~$r_L$, тому $\kappa_\perp\sim nr_L^2\nu\sim\kappa_\parallel(\nu/\omega_c)^2$~--- на багато порядків менше. Поперек
поля, до речі, за класичною оцінкою домінують іони (їхній $r_L$ більший), а~не електрони.
Тому магнітне поле є ефективною \emph{термоізоляцією} поперек силових ліній. Реальні
втрати тепла в~установках, однак, значно перевищують класичні
(див.~розд.~\ref{sec:magnetic-confinement}).

%% --------------------------------------------------------
\section{Дифузія в~плазмі. Амбіполярна дифузія}\label{sec:ambipolar}
%% --------------------------------------------------------

Нехай у~плазмі є градієнт концентрації. Електрони рухливіші за іони, тому природно було б
очікувати, що вони дифундують швидше. Але тоді виникло б розділення зарядів і~електричне
поле, яке гальмує електрони й~прискорює іони: \emph{завдяки квазінейтральності електрони
та іони дифундують разом} з~однією швидкістю. Таку дифузію називають \emphx{амбіполярною}\index{Дифузія!амбіполярна|textbf}.

Знайдемо коефіцієнт. Для $n_e=n_i=n$ потоки частинок
% ~~~~~~~~~~~~~~~~~~~~~~~~~~~~~~~~~~~~~~~~~~~~~~~~~~~~~~~~
\begin{equation*}
	\Gamma_e=-D_e\nabla n-\mu_en\vect E,\qquad
	\Gamma_i=-D_i\nabla n+\mu_in\vect E,
\end{equation*}
% ~~~~~~~~~~~~~~~~~~~~~~~~~~~~~~~~~~~~~~~~~~~~~~~~~~~~~~~~
де $D_s$~--- коефіцієнти дифузії, $\mu_s>0$~--- рухливості, пов'язані співвідношенням
Ейнштейна $D_s=\mu_sT_s/e$. Вимога $\Gamma_e=\Gamma_i\equiv\Gamma$ визначає
поле
% ~~~~~~~~~~~~~~~~~~~~~~~~~~~~~~~~~~~~~~~~~~~~~~~~~~~~~~~~
\begin{equation*}
	\vect E=-\frac{D_e-D_i}{\mu_e+\mu_i}\,\frac{\nabla n}n,
\end{equation*}
% ~~~~~~~~~~~~~~~~~~~~~~~~~~~~~~~~~~~~~~~~~~~~~~~~~~~~~~~~
а~після підстановки $\Gamma=-D_a\nabla n$ з~\emphx{амбіполярним коефіцієнтом дифузії}
% ~~~~~~~~~~~~~~~~~~~~~~~~~~~~~~~~~~~~~~~~~~~~~~~~~~~~~~~~
\begin{equation}\label{eq:ambipolar}
	\tcbhighmath{D_a=\frac{\mu_eD_i+\mu_iD_e}{\mu_e+\mu_i}=\frac{\mu_e\mu_i\,(T_e+T_i)}{e(\mu_e+\mu_i)}\approx D_i\left(1+\frac{T_e}{T_i}\right)}.
\end{equation}
% ~~~~~~~~~~~~~~~~~~~~~~~~~~~~~~~~~~~~~~~~~~~~~~~~~~~~~~~~
Останнє наближення справедливе при $\mu_e\gg\mu_i$. Отже, дифузію визначають
\emph{повільні} частинки (іони), але з~поправкою на температуру електронів: електрони
\enquote{тягнуть} іони за собою, і~$D_a$ у~$1+T_e/T_i$ разів більший за $D_i$.

Поле $\vect E\approx-(T_e/e)\,\nabla\ln n$ (при $D_e\gg D_i$) означає, що
електрони розподілені за Больцманом $n\propto\exp(e\varphi/T_e)$ в~потенціалі
$\varphi$ цього поля. Для його створення досить мізерного розділення зарядів: $\delta
n/n\sim(\lambda_D/L)^2\ll1$ (розд.~\ref{sec:plasma-basics}). Амбіполярна дифузія визначає,
наприклад, загасання розрядної плазми в~посудині після вимкнення розряду.
Формула~\eqref{eq:ambipolar} відповідає відсутності магнітного поля.

\subsection*{Дифузія поперек магнітного поля}\index{Дифузія!поперек магнітного поля|textbf}
У~замагніченій плазмі крок поперек поля дорівнює $r_L$ (кожне зіткнення зсуває ведучий центр на $\sim r_L$), а~не довжині вільного пробігу, тому
$D_\perp\sim r_L^2\nu$ (розд.~\ref{sec:particle-motion}), $D_\perp/D_\parallel\sim(\nu/\omega_c)^2$.
Запам'ятайте: \emph{у~магнітному полі зіткнення допомагають дифузії поперек $\vect B$ ($D_\perp\propto\nu$) і~заважають уздовж ($D_\parallel\propto1/\nu$)}.
У~\emph{повністю іонізованій} плазмі електрон-іонні зіткнення зберігають сумарний імпульс, і~потоки
електронів та іонів поперек поля автоматично рівні: $\nu_{ie}r_{Li}^2\sim\nu_{ei}r_{Le}^2$
(бо $\nu_{ie}=\nu_{ei}\,m_e/m_i$, а~$r_{Li}^2\sim(m_i/m_e)r_{Le}^2$ за однакових $T$). Класична поперечна дифузія є
\emph{власне амбіполярною}: додаткове електричне поле не потрібне, а~коефіцієнт
$D_\perp\sim\nu_{ei}r_{Le}^2\propto n\,T^{-1/2}/B^2$ дуже малий (для $n\approx10^{14}$~см$^{-3}$, 10~кеВ й~5~Тл
$\sim0{,}1$~см$^2$/с; для порівняння, бомівська дифузія $D_B=cT/16eB\approx1{,}3\times10^6$~см$^2$/с
більша в~$\sim10^7$ разів, див. розд.~\ref{sec:magnetic-confinement}). У~слабко іонізованому газі, де магнітяться лише електрони, повільнішими стають саме
вони, і~амбіполярна дифузія поперек поля визначається ними:
$D_{a\perp}\approx D_{e\perp}(1+T_i/T_e)$, $D_{e\perp}=D_e/(1+\omega_{ce}^2/\nu^2)$.

%% --------------------------------------------------------
\section{Як це вимірюють: температура з~опору і~з~ширини ліній}\label{sec:collision-diagnostics}
%% --------------------------------------------------------

Зіткнення залишають у~плазмі два «відбитки», з~яких можна прочитати температуру без зонда.

\emph{Температура електронів з~опору.}\index{Діагностика плазми!електронна температура з опору} Оскільки $\sigma\propto T_e^{3/2}$~\eqref{eq:spitzer}, а~густина струму $j$ і~поле $E$ в~омічному розряді вимірюються
електрично (за~струмом плазми й~напругою на обході, розд.~\ref{sec:magnetic-diagnostics}), то
\begin{equation}\label{eq:spitzer-temperature}
	\tcbhighmath{T_e[\text{еВ}]\approx\left(\frac{\sigma\ln\Lambda}{1{,}7\times10^{14}}\right)^{2/3},\qquad \sigma=\frac{j}{E}.}
\end{equation}
\begin{Estimate}
Для $I=1$~МА, $a=0{,}5$~м ($j\approx1{,}3\times10^{6}$~А/м$^2$) і~$E=0{,}1$~В/м маємо
$\sigma\approx1{,}3\times10^{7}~\text{См/м}\approx1{,}2\times10^{17}~\mathrm{с}^{-1}$, і~за $\ln\Lambda=15$ отримуємо $T_e\approx0{,}5$~кеВ.
\end{Estimate}
Метод дає усереднену по перерізу
температуру й~чутливий до~домішок: ефективний заряд $Z_{\rm eff}>1$ збільшує опір, тому $T_e$ без~поправки виходить заниженою.

\emph{Температура іонів з~допплерівського розширення.}\index{Діагностика плазми!допплерівське розширення}\index{Допплерівське розширення ліній|textbf} Іон з~масою $M$ випромінює на довжині хвилі $\lambda_0$; через тепловий рух лінія має максвеллівський профіль із~повною шириною на половині висоти
\begin{equation}\label{eq:doppler-width}
	\frac{\Delta\lambda}{\lambda_0}=\sqrt{\frac{8\ln2\,T_i}{Mc^2}}.
\end{equation}
\begin{Estimate}
Для лінії $\mathrm{H}_\alpha$ ($\lambda_0=656{,}3$~нм) за $T_i=100$~еВ маємо $\Delta\lambda/\lambda_0\approx7{,}7\times10^{-4}$, тобто $\Delta\lambda\approx0{,}5$~нм~--- легко вимірюється спектрометром.
\end{Estimate}
Важчі домішки (C, O) дають вужчі лінії, але ті самі $T_i$ за умови, що іони встигли вирівняти температуру з~основними (зіткнення, розд.~\ref{sec:plasma-collisions}).

\emph{Густина з~штарківського розширення.}\index{Діагностика плазми!штарківське розширення}\index{Штарківське розширення ліній|textbf} Мікрополя заряджених частинок розщеплюють рівні водню, і~ширина лінії зростає як $n_e^{2/3}$:
для $\mathrm{H}_\beta$ вона сягає $\approx1$~нм за $n_e\sim10^{16}$~см$^{-3}$, тож метод підходить для щільної холодної плазми (розряди, дуги), а~в~гарячій
термоядерній плазмі штарківська ширина мала на тлі допплерівської.

Обмеження: спектральна лінія несе сигнал із~усієї хорди, тому для локальних вимірів потрібне томсонівське розсіяння (розд.~\ref{sec:thomson}).

%% --------------------------------------------------------
\begin{Summary}
\begin{itemize}
\item Кулонівські зіткнення~--- накопичення багатьох малих відхилень на далеких відстанях, а~не рідкісні лобові удари.
\item Вони визначають опір плазми: провідність Спітцера $\sigma\propto T^{3/2}$, від густини вона залежить лише через $\ln\Lambda$.
\item Теплопровідність плазми переносять переважно електрони, бо вони легші й~швидші за іони.
\item Дифузія плазми амбіполярна: електричне поле зв'язує електрони та іони, вирівнюючи їхні потоки.
\item Зіткнення дають діагностику: $T_e$ з~опору (розд.~\ref{sec:collision-diagnostics}), $T_i$ з~допплерівської ширини лінії, $n_e$ зі~штарківської.
\end{itemize}
\end{Summary}

\Needspace{9\baselineskip}
\tocunnum{Задачі}
\section*{Задачі}
%% --------------------------------------------------------

\begin{problem}\label{prb:PlasmaConductivity}\lvl{2}
	Як змінюється провідність водневої плазми при зростанні температури від $10$ до
	$100$~еВ (вважайте $\ln\Lambda$ сталим)? Знайдіть $\sigma$ при $T_e=1$~кеВ, $\ln\Lambda=15$ і~порівняйте з~провідністю міді. Як змінюється
	з~температурою питома потужність омічного нагрівання при заданій густині струму?
	\begin{solution}
		$\sigma\propto T^{3/2}$: зростання в~$31{,}6$ раза; $\sigma(1~\text{кеВ})\approx3{,}6\times10^{17}~\mathrm{с}^{-1}$ (мідь: $\approx5\times10^{17}~\mathrm{с}^{-1}$); $j^2/\sigma\propto T^{-3/2}$.
	\end{solution}
\end{problem}

\begin{problem}\label{prb:PlasmaThermal}\lvl{2}
	У~скільки разів змінюються електронна теплопровідність і~провідність при нагріванні
	від $10$ до $100$~еВ? Як змінюється відношення $\kappa_e/\sigma T_e$?
	\begin{solution}
		$\kappa_e$ зростає в~316 разів, $\sigma$ в~31,6 раза; $\kappa_e/\sigma\propto T_e$ зростає в~10 разів, а~$\kappa_e/(\sigma T_e)$ не змінюється
		(аналог закону Відемана--Франца).
	\end{solution}
\end{problem}

\begin{problem}\label{prb:Dreicer}\lvl{2}
	Знайдіть поле Драйсера для $n=10^{14}$~см$^{-3}$, $T_e=10$~кеВ і~для $n=10^{13}$~см$^{-3}$, $T_e=1$~кеВ
	($\ln\Lambda=15$). Для $E=0{,}1$~В/м оцініть критичну швидкість $v_c/v_{Te}$ і~частку втікаючих електронів.
	\begin{solution}
		в~обох випадках $E_D\approx3{,}9$~В/м (відношення $n/T_e$ однакове);
	$v_c/v_{Te}=\sqrt{E_D/E}\approx6{,}2$, частка $\sim\exp(-E_D/2E)\approx3\times10^{-9}$.
	\end{solution}
\end{problem}

\begin{problem}\label{prb:PlasmaAmbipolar}\lvl{2}
	Знайдіть $D_a/D_i$ для $T_e=T_i$ та $T_e=10T_i$. Покажіть, що при $D_e\gg D_i$
	амбіполярне поле $\vect E=-(T_e/e)\nabla\ln n$, тобто електрони
	розподілені за Больцманом.
	\begin{solution}
		$D_a=2D_i$ при $T_e=T_i$ і~$D_a=11D_i$ при $T_e=10T_i$.
	\end{solution}
\end{problem}

\begin{problem}\label{prb:LnLambdaPractice}\lvl{2}
	Використовуючи табл.~\ref{tab:coulomb-log} і~формулу~\eqref{eq:nu-ei}, оцініть $\ln\Lambda$ та
	частоту $\nu_{ei}$ для трьох плазм: іоносфери ($n\sim10^{6}~\text{см}^{-3}$, $T\sim0{,}1$~еВ),
	сонячної корони ($n\sim10^{8}~\text{см}^{-3}$, $T\sim100$~еВ) і~лабораторної плазми
	($n\sim10^{13}~\text{см}^{-3}$, $T\sim10$~еВ). Яка з~них найбільш зіткненнєва щодо $\omega_p$?
	\begin{solution}
		іоносфера: $\ln\Lambda\approx12{,}6$, $\nu_{ei}\approx10^{3}~\mathrm{с}^{-1}$, $\nu_{ei}/\omega_{pe}\approx2\times10^{-5}$;
	корона: $\ln\Lambda\approx19$, $\nu_{ei}\approx6~\mathrm{с}^{-1}$, $\nu_{ei}/\omega_{pe}\approx10^{-8}$;
	лабораторна: $\ln\Lambda\approx11{,}5$, $\nu_{ei}\approx10^{7}~\mathrm{с}^{-1}$, $\nu_{ei}/\omega_{pe}\approx6\times10^{-5}$~--- найбільш зіткненнєва.
	\end{solution}
\end{problem}

\begin{problem}\label{prb:TokamakMFP}\lvl{2}
	У~токамаку $n=10^{14}~\text{см}^{-3}$, $T_e=T_i=10$~кеВ, $\ln\Lambda\approx17$, $R=3$~м.
	(а)~Оцініть $\tau_e\approx3{,}4\times10^{5}\,T_e^{3/2}[\text{еВ}]/(n\ln\Lambda)$~с, довжину вільного
	пробігу $\ell_e=v_{Te}\tau_e$ і~порівняйте її з~довжиною обходу тора $2\pi R$.
	(б)~За який час вирівнюються $T_e$ та $T_i$? Оцініть $\tau_{eq}\sim(m_i/2m_e)\,\tau_e$ і~порівняйте
	з~часом утримання енергії $\sim1$~с.
	\begin{solution}
		$\tau_e\approx2\times10^{-4}$~с, $\ell_e\approx8$~км $\gg2\pi R\approx19$~м;
	$\tau_{eq}\sim0{,}2$~с для водню.
	\end{solution}
\end{problem}

\begin{problem}\label{prb:ClassicalVsBohm}\lvl{2}
	Порівняйте класичну дифузію $D_\perp\sim r_{Le}^2\nu_{ei}$ з~бомівською $D_B=cT/(16eB)$ для
	параметрів із~задачі~\ref{prb:TokamakMFP} при $B=5$~Тл. У~скільки разів вони відрізняються
	і~що це означає для утримання?
	\begin{solution}
		$\nu_{ei}\approx5\times10^{3}~\mathrm{с}^{-1}$, $r_{Le}\approx5\times10^{-3}$~см, $D_\perp\approx0{,}1$~см$^2$/с;
	$D_B\approx1{,}3\times10^{6}$~см$^2$/с, тобто в~$\sim10^{7}$ разів більший. Якби дифузія була бомівською,
	утримання в~реакторних масштабах було б неможливим; реальний (аномальний) перенос лежить між цими межами.
	\end{solution}
\end{problem}

\begin{problem}\label{prb:SpitzerTe}\lvl{2}
	У~омічному розряді $j=2\times10^{6}$~А/м$^2$, $E=0{,}2$~В/м, $\ln\Lambda=15$. Оцініть $T_e$ за
	формулою~\eqref{eq:spitzer-temperature}. Лінія $\mathrm{H}_\alpha$ має ширину $\Delta\lambda=0{,}8$~нм; оцініть $T_i$.
	\begin{solution}
		$\sigma=10^{7}$~См/м $=9\times10^{16}~\mathrm{с}^{-1}$, $T_e\approx0{,}4$~кеВ; $T_i\approx250$~еВ
	($\Delta\lambda/\lambda_0=1{,}22\times10^{-3}$, $T_i=(\Delta\lambda/\lambda_0)^2Mc^2/8\ln2$).
	\end{solution}
\end{problem}

%% --------------------------------------------------------
%% --------------------------------------------------------

\begin{Questions}
\begin{quizlist}
\item Чому провідність повністю іонізованої плазми зростає з~температурою й~не
		  залежить від концентрації?

\item Чому в~слабкозв'язаній плазмі ($N_D\gg1$) $\nu_{ei}\ll\omega_p$, і~чому це ще не означає, що
		  довільний процес беззіткнювальний? Яку роль грає кулонівський логарифм?

\item Що таке амбіполярна дифузія? Чому $D_a\approx D_i(1+T_e/T_i)$ і~яка роль
		  квазінейтральності?

\item Що таке поле Драйсера? Чому сила тертя на швидкий електрон спадає зі~швидкістю?

\item Чому класична дифузія поперек поля у~повністю іонізованій плазмі автоматично амбіполярна?

\item Чому теплопровідність плазми $\propto T^{5/2}$ і~не залежить від
		  концентрації? Чому вона різна вздовж і~поперек поля?

\item Чому зі~зростанням температури плазма стає кращим провідником, а~метал~--- гіршим?

\item Чому зіткнення сприяють дифузії поперек поля, але гальмують дифузію вздовж нього?

\item Як за~напругою на обході та струмом оцінити $T_e$ омічної плазми? Чому результат залежить від~$Z_{\rm eff}$?

\end{quizlist}
\end{Questions}

\Closesolutionfile{answer}
